\documentclass[conference]{IEEEtran}
\IEEEoverridecommandlockouts

% The preceding line is only needed to identify funding in the first footnote. If that is unneeded, please comment it out.
\usepackage{cite}
\usepackage{amsmath,amssymb,amsfonts}
\usepackage{algorithmic}
\usepackage{graphicx}
\usepackage{textcomp}
\usepackage{xcolor}
\usepackage{multirow}
\usepackage{amsthm}
\usepackage{mathrsfs}
\usepackage{manyfoot}
\usepackage{booktabs}
\usepackage{algorithm}
\usepackage{listings}
\usepackage{times}
\usepackage{epsfig}
\usepackage[export]{adjustbox}
\usepackage{tabularx}

\def\BibTeX{{\rm B\kern-.05em{\sc i\kern-.025em b}\kern-.08em
    T\kern-.1667em\lower.7ex\hbox{E}\kern-.125emX}}

\begin{document}

\title{Super Intelligence and General Intelligence\\
%{\footnotesize \textsuperscript{*}Note: Sub-titles are not captured in Xplore and
%should not be used}
%\thanks{Identify applicable funding agency here. If none, delete this.}
}

\author{
\IEEEauthorblockN{Zanyar Zohourianshahzadi}
\IEEEauthorblockA{
\textit{Artificial Intelligence, Computer Science, and CyberSecurity} \\
\textit{B.S. CS \& IT: Azad University of Tehran South Branch} \\
\textit{M.S. IT \& CyberSecurity: Colorado Technical University} \\
\textit{Ph.D. AI \& CS: University of Colorado Colorado Springs} \\
Colorado Springs, Colorado, USA \\
Email: {zanyarz@me.com \& zanyarz@gmail.com}
}
}

\maketitle

\begin{abstract}
Artificial general intelligence (AGI) and superintelligence are frequently treated as successive points on a single scale of cognitive performance. This paper argues that they represent distinct capabilities. Building on intelligence as information-dependent decision capability and AGI as integrated multimodal agency, we define general intelligence by an agent’s capacity to coordinate perception, reasoning, memory, tools, decisions, and actions across domains. Superintelligence requires an additional epistemic transition: the ability to expand the frontier of knowledge beyond both what humanity knows and the questions humanity has already learned to ask. We distinguish three related processes. Knowledge fusion establishes consequential relations among previously separated bodies of knowledge. Knowledge discovery identifies previously unknown phenomena, mechanisms, regularities, or formal objects. Knowledge invention occurs when a machine originates a novel epistemic structure that becomes consequential, independently validated, and incorporable into collective knowledge. We formalize this distinction through the Knowledge-Invention Proposition and propose a staged evaluation framework based on frontier verification, unknown-unknown origination, causal attribution, independent validation, consequentiality, and repeatability. The framework distinguishes machine assistance, machine discovery, and validated machine knowledge invention while addressing contamination, human specification, epistemic opacity, and dual-use risk. Finally, we maintain that consciousness, general intelligence, and superintelligence are analytically independent. Superintelligence is therefore not defined by the quantity of knowledge stored or performance on known tasks, but by the repeatable capacity of a generally intelligent machine to originate validated knowledge beyond the human epistemic frontier.
\end{abstract}

\begin{IEEEkeywords}
Super Intelligence, AGI, General Intelligence, Artificial Consciousness, Machine Consciousness
\end{IEEEkeywords}


\section{\textbf{Introduction}}
\label{sec:introduction}

Intelligence, consciousness, general intelligence, and superintelligence are
frequently placed on a single ascending scale. This treatment obscures the
different questions addressed by these concepts. Intelligence concerns a
system's capacity to use information in selecting among possible actions,
states, or outputs. Consciousness concerns the organization and, in its
phenomenal sense, the subjective character of experience. General intelligence
concerns the breadth, transferability, and adaptive organization of agency.
Superintelligence, we argue, concerns a further capability: moving beyond the
collective human frontier by inventing validated knowledge that humanity did
not previously possess or recognize as an object of inquiry.

Our earlier work defined intelligence at its minimal foundation as
information-dependent selection while distinguishing intelligence from both
functional and phenomenal consciousness \cite{zohourian2026intelligence}. On
this account, intelligence is graded rather than binary. A simple conditional
selection may instantiate an elementary decision structure, whereas advanced
systems can coordinate memory, inference, planning, metacognition, and action.
Increasing sophistication can therefore produce more organized intelligence
without proving the existence or character of subjective experience.

Consciousness should likewise not be restricted in advance to one biological,
humanlike, or computational form. We previously argued that phenomenal,
observable, and functional claims about consciousness must be separated and
that language can serve as a unifying interface among specialized artificial
minds \cite{zohourian2026language}. Language is not merely an input--output
modality under this account. It can provide a shared representational medium
through which perceptual, reasoning, memory, and supervisory systems exchange
information. Such coordination may enable increasingly integrated levels of
functional artificial consciousness without, by itself, establishing
phenomenal experience.

These distinctions become essential when considering artificial general
intelligence (AGI). Our account of AGI as integrated multimodal agency treats
general intelligence and consciousness as related but non-identical dimensions
\cite{zohourian2026agi}. Neither multimodality nor consciousness alone supplies
general intelligence. AGI requires the persistent coordination of perception,
memory, reasoning, objectives, planning, metacognition, and adaptive action
across unfamiliar tasks, domains, and environments. This view is consistent
with accounts that relate intelligence to an agent's ability to achieve goals
across environments rather than to performance on one fixed task
\cite{legg2007universal} and with broader treatments of AGI as adaptable,
domain-general intelligence rather than a collection of isolated task systems
\cite{goertzel2014agi}.

Future work will introduce \emph{AGI-Former}, a neural architecture intended to
operationalize this conception of integrated multimodal agency. The present
paper does not specify that architecture. It instead addresses the conceptual
problem that follows from it: even if a machine realizes AGI, what additional
capability would make it superintelligent?

\subsection{\textbf{From General Intelligence to Superintelligence}}

General intelligence and superintelligence are sometimes treated as
interchangeable descriptions of an increasingly capable system. They instead
identify different functional achievements. AGI denotes integrated,
transferable, and adaptive agency across tasks, domains, modalities, and
environments. Superintelligence denotes capability beyond the collective human
frontier: the capacity of a generally intelligent machine to originate
validated knowledge that humanity did not previously possess or recognize as
an object of inquiry. This retains the established comparative meaning of
superintelligence as capability exceeding human intelligence
\cite{bostrom2014superintelligence}, while specifying the epistemic capability
that distinguishes it from AGI in the present framework.

An AGI could therefore remain bounded by existing human knowledge. It might
reason across disciplines, operate tools, adapt plans, and solve difficult
problems while deriving its objectives and conceptual boundaries from questions
already posed by humans. These abilities would establish general agency, but
they would not alone demonstrate superintelligence. The transition occurs when
the machine ceases to be only a general user of known science and becomes an
originator of new science.

This distinction can be expressed directly:

\begin{equation}
\begin{aligned}
\operatorname{AGI}(x)
&\not\Rightarrow \operatorname{SI}(x),\\
\operatorname{SI}(x)
&\Rightarrow \operatorname{AGI}(x).
\end{aligned}
\label{eq:agi_si_relation}
\end{equation}

The additional condition is not greater speed, scale, memory, or performance
within a predetermined task. It is machine-originated knowledge invention. To
state this condition precisely, the distinction between known unknowns and
unknown unknowns must first be established.

\subsection{\textbf{Known Unknowns and Unknown Unknowns}}

Scientific research commonly begins with a \emph{known unknown}: a question,
target, mechanism, or missing value that researchers have already identified
but not yet resolved. A machine may answer such a question more quickly or
accurately than human researchers and thereby make a major scientific
contribution. The machine is operating beyond prior human performance, but the
epistemic destination was already represented within human inquiry.

An \emph{unknown unknown} is stronger. It is a consequential question,
relation, phenomenon, explanatory possibility, or field of inquiry that
humanity has neither resolved nor recognized as a target. Unknown unknowns are
not simply entries missing from a database. Before they are exposed, humans do
not know that a particular question should be asked, that two bodies of
knowledge should be connected, that an observed anomaly signifies a new
mechanism, or that a new conceptual object is possible.

Let $\mathcal{K}_{H,t}$ denote the collective body of justified human knowledge
at time $t$, and let $\mathcal{Q}_{H,t}$ denote the set of questions, hypotheses,
and research targets explicitly recognized by the human scientific community.
For a candidate epistemic contribution $\kappa$,

\begin{equation}
\begin{aligned}
\kappa \in \mathcal{Q}_{H,t}
\land \kappa \notin \mathcal{K}_{H,t}
&\quad\Longrightarrow\quad
\kappa\ \text{is a known unknown},\\
\kappa \notin \mathcal{Q}_{H,t}
\land \kappa \notin \mathcal{K}_{H,t}
&\quad\Longrightarrow\quad
\kappa\ \text{is an unknown unknown}.
\end{aligned}
\label{eq:known_unknowns}
\end{equation}

The second relation cannot be established solely by observing that a generated
sentence is absent from a training corpus. Semantic novelty, scientific
significance, and epistemic validity are different properties. A system must
identify or construct the previously unrecognized object of inquiry, explain
its relation to existing knowledge, and produce evidence through which the
claim can be independently assessed.

Existing AI-for-science results demonstrate important progress without yet
collapsing this distinction. AI has recovered symbolic relations from data
\cite{udrescu2020aifeynman}, guided mathematicians toward new conjectures and
theorems \cite{davies2021mathematics}, predicted protein structures at high
accuracy \cite{jumper2021alphafold}, and expanded computational catalogues of
candidate stable materials \cite{merchant2023materials}. These achievements
show that machine learning can contribute to discovery at the scientific
frontier. They also remain embedded in research problems, objectives, data,
and validation practices established by human investigators. Under the present
framework, their scientific importance is not in question; what remains to be
established is whether a machine can itself expose an unknown unknown and
originate the knowledge through which it becomes understood.

\subsection{\textbf{Knowledge Invention}}

We use \emph{knowledge invention} for the machine-originated production of a
new epistemic object that lies beyond the collective human frontier and becomes
justified through appropriate validation. The term includes the invention of
new explanatory theories, conceptual structures, scientific methods,
architectures, or research directions. It does not imply that natural laws are
arbitrarily created. A machine may discover a previously unknown phenomenon or
relation, but the representation, explanation, and generalizable body of
knowledge constructed from that discovery constitute an invented epistemic
contribution.

Knowledge fusion and knowledge discovery are principal routes toward knowledge
invention. Knowledge fusion connects bodies of evidence, concepts, models, or
disciplines whose consequential relation was not previously recognized.
Knowledge discovery identifies a previously unknown phenomenon, mechanism,
regularity, entity, or law. Either process may produce knowledge invention when
the machine converts its result into a novel, consequential, and validated
addition to collective knowledge:

\begin{equation}
\begin{aligned}
\operatorname{Fusion}_{x}
&\longrightarrow \kappa_{F},\\
\operatorname{Discovery}_{x}
&\longrightarrow \kappa_{D},\\
\operatorname{Validate}(\kappa_{F}\lor\kappa_{D})=1
&\longrightarrow \operatorname{KI}_{x}(\kappa),
\end{aligned}
\label{eq:routes_to_invention}
\end{equation}

where $\operatorname{KI}_{x}(\kappa)$ denotes knowledge invention attributable
to machine $x$. The arrows represent possible epistemic pathways rather than a
claim that every fusion or discovery yields invention. Most recombinations are
not consequential, many hypotheses are false, and some discoveries answer
questions already posed by humans. Knowledge invention therefore requires a
stronger conjunction:

\begin{equation}
\begin{aligned}
\operatorname{KI}_{x}(\kappa)
\Longleftrightarrow{}&
\kappa \notin \mathcal{K}_{H,t}\\
&\land \kappa \notin \mathcal{Q}_{H,t}\\
&\land \operatorname{Originate}_{x}(\kappa)\\
&\land \operatorname{Consequential}(\kappa)\\
&\land \operatorname{Validate}(\kappa)=1.
\end{aligned}
\label{eq:knowledge_invention}
\end{equation}

The origination condition does not require isolation from humans. Scientific
work is commonly distributed across people, instruments, institutions, and
computational systems. It requires that the machine make the causally essential
move that exposes or constructs the unknown unknown, rather than merely execute
a fully specified human procedure. The consequentiality condition excludes
arbitrary novelty. The validation condition may be satisfied through
controlled experiment, mathematical proof, reproducible computation,
successful engineering realization, or another method appropriate to the
domain. Until validation occurs, the output remains a candidate invention or
hypothesis rather than knowledge.

Recent systems already integrate language models with literature search, code,
planning, laboratory interfaces, and experimental execution
\cite{boiko2023autonomous}, while autonomous laboratories combine computation,
machine learning, robotics, and empirical feedback in closed experimental loops
\cite{szymanski2023autonomous}. Reviews of AI-enabled science similarly show
that learning systems now participate across hypothesis generation, prediction,
experimental design, and analysis \cite{wang2023scientific}. These systems
provide architectural and empirical precursors to machine knowledge invention.
They do not by their existence satisfy Equation~\eqref{eq:knowledge_invention},
because scientific autonomy within a human-specified objective is not identical
to originating an unrecognized object of inquiry.

\subsection{\textbf{The Defining Transition to Superintelligence}}

Under the framework proposed here, AGI is necessary but not sufficient for
superintelligence. An AGI may transfer knowledge, adapt to unfamiliar
environments, and pursue objectives across domains while remaining bounded by
the human knowledge and research targets available to it. Superintelligence
begins when general intelligence can originate and validate knowledge beyond
that collective frontier:

\begin{equation}
\begin{aligned}
\operatorname{SI}(x)
\Longleftrightarrow{}&
\operatorname{AGI}(x)\\
&\land \exists\kappa\;\operatorname{KI}_{x}(\kappa)\\
&\land
\mathcal{K}_{H,t+1}
=\mathcal{K}_{H,t}\cup\{\kappa\},\\
\operatorname{SI}(x)
&\Rightarrow \operatorname{AGI}(x),\\
\operatorname{AGI}(x)
&\not\Rightarrow \operatorname{SI}(x).
\end{aligned}
\label{eq:superintelligence_definition}
\end{equation}

The central thesis of this paper is therefore that \emph{superintelligence is
general intelligence capable of originating validated knowledge beyond the
collective human frontier, including the identification and resolution of
unknown unknowns}. Knowledge fusion and knowledge discovery are not three
parallel requirements alongside knowledge invention. They are principal
generative routes through which knowledge invention may occur.

This definition does not equate superintelligence with omniscience,
infallibility, consciousness, model size, or universal superiority on every
task. It makes a narrower and testable claim: the defining evidence for
superintelligence is a machine-originated epistemic contribution that humanity
did not previously possess or explicitly seek, that materially extends human
understanding, and that survives independent validation. The remainder of the
paper distinguishes this capability from general intelligence,
examines knowledge fusion and discovery as routes to invention, and develops
criteria for attributing knowledge invention to an artificial system without
mistaking recombination, hallucination, or unverified novelty for science.

\section{\textbf{From General Intelligence to Superintelligence}}
\label{sec:agi_to_si}

Section~\ref{sec:introduction} proposed that superintelligence begins when a
generally intelligent machine originates validated knowledge beyond the
collective human frontier. This section develops that distinction more
precisely. The transition from AGI to superintelligence is not defined by a
larger model, a longer context window, access to more tools, or improved
performance within an existing task. It is defined by a change in the machine's
epistemic role: from a system that operates across established knowledge and
human-specified problems to a system that can expose, formulate, and resolve
problems that humanity had not recognized.

\subsection{\textbf{The Functional Boundary of General Intelligence}}

General intelligence concerns the organization of agency. An AGI must transfer
knowledge across domains, recompose existing capabilities for unfamiliar
problems, maintain goals over time, revise plans after failure, and coordinate
perception, memory, reasoning, metacognition, and action
\cite{zohourian2026agi}. Formal definitions of intelligence likewise emphasize
the capacity to achieve goals across a wide range of environments
\cite{legg2007universal}, while broader AGI accounts distinguish general
adaptability from competence on isolated tasks \cite{goertzel2014agi}.

Let $\mathcal{E}$ denote environments, $\mathcal{T}$ tasks, and $\mathcal{M}$
modalities. For agent $x$, let $T_x$, $R_x$, and $A_x$ denote transfer,
recomposition, and adaptation. The functional scope of general agency can be
represented as

\begin{equation}
\begin{aligned}
\mathcal{G}(x)
&= \Phi_x(T_x,R_x,A_x),\\
\Phi_x:
\mathcal{E}\times\mathcal{T}\times\mathcal{M}
&\longrightarrow
\text{coherent goal-directed behavior},
\end{aligned}
\label{eq:general_agency}
\end{equation}

where $\Phi_x$ denotes the persistent organization through which these
capacities inform and constrain one another. No single component is sufficient.
Transfer without adaptation may fail when prior assumptions break;
recomposition without coherent goals may produce locally useful but globally
inconsistent actions; and broad task exposure may reflect memorized coverage
rather than general agency.

Generality also cannot be inferred from the mere presence of many components.
A framework may contain language, vision, memory, search, code execution, and
robotic interfaces while relying on a human to select tools, reconcile their
outputs, detect failure, and preserve the objective. Framework capacity,
realized agency, and observed performance must therefore remain distinct:

\begin{equation}
\begin{aligned}
\operatorname{Cap}(F)
&\not\equiv \operatorname{Realize}(x,F),\\
\operatorname{Realize}(x,F)
&\not\equiv \operatorname{Perf}(x,\tau),
\end{aligned}
\label{eq:capacity_realization}
\end{equation}

where $F$ is a computational framework and $\tau$ is an evaluated task. An AGI
realizes coordinated general agency; it need not be omniscient, infallible, or
superior on every task. Its defining property is the ability to reorganize its
available capacities in pursuit of objectives across sufficiently varied and
unfamiliar conditions.

\subsection{\textbf{The Epistemic Boundary of General Intelligence}}

Nothing in the preceding definition requires AGI to create knowledge beyond
humanity. A generally intelligent machine could possess extensive human
knowledge, apply it across disciplines, generate plans, conduct analyses, and
solve previously unsolved but already formulated problems. Such a system could
be transformative while remaining epistemically bounded by the concepts,
questions, and objectives inherited from human inquiry.
Reviews of AI-enabled scientific discovery already distinguish multiple roles
for machines across hypothesis generation, modelling, experimental design, and
analysis \cite{wang2023scientific}. That breadth of participation motivates a
more precise question of attribution: whether the machine merely operates
within the human research agenda or originates an epistemic object beyond it.

Let $\mathcal{K}_{H,t}$ be collective justified human knowledge at time $t$,
and let $\mathcal{Q}_{H,t}$ be the set of questions, hypotheses, and research
targets recognized by humans. An AGI operating within the human epistemic
boundary may be represented by

\begin{equation}
\begin{aligned}
x:
(\mathcal{K}_{H,t},\mathcal{Q}_{H,t},\mathcal{O}_{t})
&\longrightarrow
(\mathcal{Y}_{t},\mathcal{A}_{t}),\\
\mathcal{Q}_{x,t}
&\subseteq \mathcal{Q}_{H,t},
\end{aligned}
\label{eq:agi_epistemic_boundary}
\end{equation}

where $\mathcal{O}_{t}$ denotes available observations,
$\mathcal{Y}_{t}$ generated conclusions or predictions, and
$\mathcal{A}_{t}$ actions. The subset relation expresses the limiting case in
which the machine's research agenda remains contained within questions humans
have already recognized. The machine may produce an answer absent from
$\mathcal{K}_{H,t}$ and thereby resolve a known unknown, but it does not yet
originate the object of inquiry itself.

This distinction is visible in current AI-assisted science. Machine learning
has helped recover symbolic laws from experimental data
\cite{schmidt2009laws,udrescu2020aifeynman}, guided
mathematicians toward conjectures and theorems \cite{davies2021mathematics},
predicted protein structures \cite{jumper2021alphafold}, and proposed large
numbers of candidate stable materials \cite{merchant2023materials}. Language
models connected to tools and laboratory automation can design, plan, and
execute chemistry experiments \cite{boiko2023autonomous}, while autonomous
laboratories can perform closed-loop materials synthesis
\cite{szymanski2023autonomous}. Earlier robot-scientist systems likewise
generated and experimentally tested functional-genomics hypotheses
\cite{king2004robot} and demonstrated end-to-end automation of iterative
scientific investigation \cite{king2009automation}. These are consequential advances at the
scientific frontier. They also begin from domains, objectives, datasets,
experimental systems, and validation conditions established by human
researchers. Their results demonstrate increasingly powerful scientific
instruments and agents, not by themselves the invention of an unknown unknown.

\subsection{\textbf{Autonomous Problem Formation}}

The first additional requirement for superintelligence is autonomous problem
formation. The machine must be capable not only of answering a question but of
identifying that a consequential question exists. This may occur when the
system detects an unexplained regularity, recognizes a contradiction across
disciplines, constructs a new conceptual variable, identifies an assumption
shared by existing theories, or notices that available observations support a
previously unrepresented phenomenon. Work on automated identification of
latent state variables from high-dimensional observations illustrates progress
toward machines constructing representations that were not specified directly
by researchers \cite{chen2022variables}. Such representation discovery is an
important precursor, although it does not alone establish autonomous problem
formation in the stronger sense proposed here.

Let $q_{x}$ be a machine-originated research question. Autonomous problem
formation requires

\begin{equation}
\begin{aligned}
q_x \notin \mathcal{Q}_{H,t},
\qquad
q_x \notin \mathcal{K}_{H,t},\\
\operatorname{Ground}(q_x,\mathcal{K}_{H,t},\mathcal{O}_{t})=1,
\qquad
\operatorname{Consequential}(q_x)=1.
\end{aligned}
\label{eq:problem_formation}
\end{equation}

The grounding condition prevents arbitrary question generation from being
mistaken for intelligence. A language model can generate indefinitely many
unusual questions through stochastic variation. Most will be trivial,
ill-posed, physically impossible, semantically incoherent, or disconnected
from evidence. A superintelligent problem formulation must arise from a
defensible relation among observations, established knowledge, unresolved
tensions, or formally represented possibilities. The consequentiality
condition further requires that resolving the question would materially alter
understanding, explanation, prediction, or technological capability.

Autonomous problem formation is therefore more demanding than novelty. It
requires epistemic judgment: the system must distinguish an unasked question
that matters from an unasked question that is merely possible. This capacity
is the point at which an unknown unknown becomes an identified object of
inquiry.

\subsection{\textbf{Causal Origination of Knowledge}}

Scientific results are commonly produced by distributed systems containing
researchers, datasets, instruments, software, models, and institutions.
Requiring a superintelligent machine to work without human participation would
therefore impose a standard not applied to human scientists. Nevertheless,
attribution cannot be granted merely because AI appears somewhere in the
workflow. The machine must make the causally essential epistemic contribution.

Let $W$ be a scientific workflow and let $\kappa$ be its resulting knowledge
claim. We define machine origination counterfactually:

\begin{equation}
\begin{aligned}
\operatorname{Originate}_{x}(\kappa\mid W)=1
\Longleftrightarrow{}&
\operatorname{Produce}(W,\kappa)=1\\
&\land
\operatorname{Produce}(W\setminus x,\kappa)=0\\
&\land
\operatorname{Specified}_{H}(\kappa)=0.
\end{aligned}
\label{eq:causal_origination}
\end{equation}

The first condition requires that the workflow produce the contribution. The
second requires that removing the machine's relevant epistemic operation would
prevent that contribution from being produced through the same workflow. The
third excludes cases in which humans fully specified the conclusion and used
the machine only to calculate, retrieve, format, or experimentally execute it.
This is a conceptual causal test rather than a demand for literal repetition of
every discovery process.

Machine origination can occur through knowledge fusion or knowledge discovery.
In fusion, the essential contribution may be recognizing that previously
separated theories, data, or disciplines instantiate a common structure. In
discovery, it may be detecting a new entity, mechanism, regularity, or law.
Neither pathway is sufficient until its output is converted into an epistemic
object that can be understood, assessed, and incorporated into knowledge.
Prior systems for symbolic-law recovery and automated hypothesis testing
demonstrate that portions of this cycle can be computationally realized
\cite{schmidt2009laws,king2004robot}; the unresolved issue is whether the
machine can originate the consequential object of inquiry rather than search
within a human-defined hypothesis space.

\subsection{\textbf{The Transition Criterion}}

The transition from AGI to superintelligence requires a complete epistemic
cycle. The machine must form or expose the previously unrecognized problem,
originate a candidate explanation or solution, connect that contribution to
existing knowledge, and enable domain-appropriate validation. Let
$\kappa_x$ denote the resulting machine-originated contribution. Then

\begin{equation}
\begin{aligned}
\operatorname{KI}_{x}(\kappa_x)
\Longleftrightarrow{}&
\operatorname{FormProblem}_{x}(q_x)\\
&\land \operatorname{Originate}_{x}(\kappa_x\mid W)\\
&\land \operatorname{Explain}_{x}
(\kappa_x,\mathcal{K}_{H,t})\\
&\land \operatorname{Validate}(\kappa_x)=1\\
&\land \operatorname{Consequential}(\kappa_x)=1.
\end{aligned}
\label{eq:knowledge_invention_cycle}
\end{equation}

Knowledge invention is not established at the moment a machine produces a
novel output. Before validation, the output is a candidate hypothesis, theory,
method, or invention. Validation may require controlled experiment, formal
proof, reproducible computation, engineering realization, convergent evidence,
or another procedure accepted by the relevant discipline. Independent human or
machine verification may participate in this process, but the validation
standard must not be defined solely by the originating system. This insistence
on domain-appropriate validation follows the broader AI-for-science literature,
in which prediction, hypothesis generation, experimentation, and analysis form
an interdependent discovery process rather than interchangeable evidence of
knowledge \cite{wang2023scientific}.

The relationship between AGI and superintelligence can consequently be stated
as

\begin{equation}
\begin{aligned}
\operatorname{SI}(x)
\Longleftrightarrow{}&
\operatorname{AGI}(x)\\
&\land \exists q_x,\kappa_x:
q_x\notin\mathcal{Q}_{H,t}\\
&\land \operatorname{KI}_{x}(\kappa_x)\\
&\land
\mathcal{K}_{H,t+1}
=\mathcal{K}_{H,t}\cup\{\kappa_x\}.
\end{aligned}
\label{eq:si_transition}
\end{equation}

Equation~\eqref{eq:si_transition} preserves the established meaning of
superintelligence as capability beyond humans \cite{bostrom2014superintelligence}
while replacing an unspecified performance comparison with an epistemic
criterion. The machine goes beyond humans not merely by doing a known task
better, but by making a validated addition to collective knowledge whose object
of inquiry humanity had not recognized.

This criterion also clarifies why the present framework does not define a
separate category of domain superintelligence. A specialized scientific model
may be indispensable to a discovery, but specialization does not establish
general intelligence. Under the terminology developed here,
superintelligence is a mode of general intelligence. Specialized systems are
scientific instruments or components of a discovery system unless they also
realize the general agency and machine-originated knowledge invention required
by Equation~\eqref{eq:si_transition}.

Knowledge fusion and knowledge discovery now require separate analysis because
they describe the principal routes by which a machine may cross the human
epistemic boundary. The following sections examine how each route can produce
knowledge invention, how their outputs differ from retrieval and recombination,
and what evidence would justify attributing the resulting knowledge to the
machine.

\section{\textbf{Knowledge Fusion as a Pathway to Knowledge Invention}}
\label{sec:knowledge_fusion}

Knowledge is distributed across disciplines, institutions, datasets,
terminologies, and historical periods. Two scientific communities may study
related structures without sharing concepts, cite separate literatures while
describing interacting mechanisms, or possess results that become
consequential only when interpreted together. This fragmentation creates a
space in which the components of a new insight may already be publicly
available even though the insight itself is absent from collective knowledge.
Knowledge fusion is the process by which an intelligent system identifies and
constructs such a missing relation.

The possibility of deriving new hypotheses from disconnected literatures has
a substantial history. Swanson described \emph{undiscovered public knowledge}
through the example of separately published biomedical findings whose latent
connection suggested a previously unrecognized hypothesis
\cite{swanson1986fish}. Later work on literature-based discovery generalized
this idea into computational methods for identifying implicit relations across
scientific corpora. Studies of scientific innovation likewise show that
atypical combinations of prior work are associated with unusually high impact
\cite{uzzi2013atypical}, while research that bridges distant regions of a
knowledge network is riskier but can yield greater rewards
\cite{foster2015tradition}. These findings motivate knowledge fusion, but the
definition developed here imposes requirements beyond bibliographic novelty.

\subsection{\textbf{From Combination to Fusion}}

Let $\mathcal{K}_{1},\ldots,\mathcal{K}_{n}$ denote bodies of justified
knowledge associated with distinct domains, communities, or representational
systems. Their aggregation is the union

\begin{equation}
\begin{aligned}
\mathcal{K}_{\cup}
&=\bigcup_{i=1}^{n}\mathcal{K}_{i}.
\end{aligned}
\label{eq:knowledge_union}
\end{equation}

Aggregation increases availability but does not create a new epistemic
relation. A search engine, database, or language model may retrieve elements
from several domains and place them in one response without changing what is
known. Knowledge fusion requires an additional contribution $\kappa_F$ that is
not explicitly contained in any source domain or in their simple union:

\begin{equation}
\begin{aligned}
\mathcal{F}_{x}
(\mathcal{K}_{1},\ldots,\mathcal{K}_{n})
&\longrightarrow \kappa_F,\\
\kappa_F
&\notin \mathcal{K}_{i}
\quad \forall i,\\
\kappa_F
&\notin \mathcal{K}_{\cup}.
\end{aligned}
\label{eq:fusion_novel_relation}
\end{equation}

Here, $\mathcal{F}_{x}$ is the fusion process realized by machine $x$, and
$\kappa_F$ is a candidate relation, explanation, hypothesis, or conceptual
structure. The final condition means that $\kappa_F$ cannot be recovered as an
explicit claim already present in the aggregated sources. Its components may
be present, but their consequential connection is not.

This distinction separates five increasingly demanding operations. Retrieval
locates an existing item. Aggregation places items together. Synthesis organizes
existing claims into a coherent account. Recombination rearranges known
elements into a configuration that may be novel in form. Fusion constructs a
relation that changes what can be explained, predicted, tested, or created.
Only the last operation is a candidate pathway to knowledge invention.

\subsection{\textbf{Latent Relations Across Knowledge Domains}}

A latent relation exists when justified claims in separated knowledge domains
jointly support an implication not represented within any participating
domain. The simplest structure is the open-discovery pattern

\begin{equation}
\begin{aligned}
A &\longrightarrow B,\\
B &\longrightarrow C,\\
A &\overset{?}{\longrightarrow} C,
\end{aligned}
\label{eq:abc_fusion}
\end{equation}

where the solid relations are established and the dashed relation is an
unrecognized candidate inferred through the shared intermediary $B$. This
structure captures the foundation of classical literature-based discovery
\cite{swanson1986fish}, but knowledge fusion need not be limited to a
three-node chain. It may connect causal graphs, mathematical structures,
experimental anomalies, measurement techniques, or explanatory models across
many domains.

The importance of a fusion does not follow from distance alone. Two remote
concepts may be unrelated, and two adjacent concepts may support a profound
inference. Empirical studies of scientific production indicate that successful
high-impact work often balances conventional foundations with atypical
combinations \cite{uzzi2013atypical}. A fusion system must therefore preserve
enough established structure to support inference while exploring relations
that prevailing disciplinary organization makes unlikely to be examined.

Let $d(\mathcal{K}_{i},\mathcal{K}_{j})$ measure conceptual or bibliographic
distance, and let $J(\kappa_F)$ represent the justification available for a
candidate fusion. Then neither novelty nor distance alone is sufficient:

\begin{equation}
\begin{aligned}
d(\mathcal{K}_{i},\mathcal{K}_{j})\gg 0
&\not\Rightarrow J(\kappa_F)>0,\\
J(\kappa_F)>0
&\not\Rightarrow
\operatorname{Consequential}(\kappa_F)=1.
\end{aligned}
\label{eq:distance_not_fusion}
\end{equation}

The first relation rejects arbitrary cross-domain analogy. The second rejects
valid but trivial connections. Knowledge fusion requires both defensible
support and consequential epistemic change.

\subsection{\textbf{Fusion as the Exposure of an Unknown Unknown}}

Knowledge fusion becomes relevant to superintelligence when it reveals not
only a new answer but a previously unrecognized object of inquiry. Suppose
$\mathcal{K}_{a}$ and $\mathcal{K}_{b}$ contain individually coherent models,
yet their joint implications produce a contradiction, unexplained regularity,
or missing causal variable. A machine may use that tension to formulate a new
question $q_F$:

\begin{equation}
\begin{aligned}
\mathcal{F}_{x}(\mathcal{K}_{a},\mathcal{K}_{b})
&\longrightarrow q_F,\\
q_F
&\notin \mathcal{Q}_{H,t},\\
\operatorname{Ground}(q_F,
\mathcal{K}_{a}\cup\mathcal{K}_{b})
&=1.
\end{aligned}
\label{eq:fusion_unknown_unknown}
\end{equation}

The question is an unknown unknown before fusion because humanity has not
represented it within the recognized research agenda. Fusion changes that
state. It exposes the question, provides its initial justification, and creates
a target for discovery or theory construction. Computational work that
identifies previously unspecified state variables from observations provides a
partial example of machines constructing new representational candidates
\cite{chen2022variables}. The stronger fusion criterion additionally requires
that the machine recognize the cross-domain significance of the constructed
object.

Scientific search is shaped by human incentives and by the organization of
existing knowledge. Network analyses show that researchers preferentially
pursue established relationships even when riskier, more innovative searches
have greater potential value \cite{foster2015tradition}. Computational systems
can search differently. Models of collective discovery suggest that strategic
selection of experiments can accelerate the exploration of a scientific
knowledge network \cite{rzhetsky2015experiments}. A fusion-oriented machine
could extend this principle by selecting not only experiments but also distant
knowledge regions whose joint structure merits investigation.

\subsection{\textbf{Fusion-Driven Knowledge Invention}}

A candidate fusion becomes knowledge invention only after it is developed into
a consequential epistemic contribution and validated. Let
$\operatorname{FKI}_{x}(\kappa_F)$ denote fusion-driven knowledge invention by
machine $x$. Then

\begin{equation}
\begin{aligned}
\operatorname{FKI}_{x}(\kappa_F)
\Longleftrightarrow{}&
\operatorname{Fuse}_{x}
(\mathcal{K}_{1},\ldots,\mathcal{K}_{n})\\
&\land \kappa_F\notin\mathcal{K}_{H,t}\\
&\land
\operatorname{Originate}_{x}(\kappa_F\mid W)\\
&\land \operatorname{Consequential}(\kappa_F)\\
&\land \operatorname{Validate}(\kappa_F)=1.
\end{aligned}
\label{eq:fusion_knowledge_invention}
\end{equation}

Equation~\eqref{eq:fusion_knowledge_invention} excludes a system that merely
retrieves a previously published interdisciplinary connection. It also excludes
a system that proposes a novel connection without evidence, repeats a human
specified synthesis, or generates an interesting analogy with no explanatory
or predictive consequence. The machine must originate the essential relation
and that relation must survive independent, domain-appropriate evaluation.

Validation depends on the type of fused contribution. A causal hypothesis may
require controlled intervention; a mathematical bridge may require proof; a
unified explanatory model may require superior prediction across the source
domains; and a proposed technological principle may require successful
realization. AI-enabled science already distributes hypothesis generation,
modelling, experimental planning, and analysis across interacting systems
\cite{wang2023scientific}. Fusion-driven knowledge invention adds the
requirement that the cross-domain relation itself originate from the machine
rather than be specified by the human workflow.

\subsection{\textbf{Failure Modes and Evidentiary Safeguards}}

Knowledge fusion is especially vulnerable to outputs that appear profound
because they connect distant concepts. At least six failure modes must be
distinguished from genuine fusion. First, \emph{lexical overlap} mistakes shared
terminology for shared mechanism. Second, \emph{analogy inflation} treats a
structural resemblance as evidence of causal identity. Third,
\emph{correlation substitution} presents statistical association as
explanation. Fourth, \emph{ontology mismatch} combines concepts whose meanings
or measurement assumptions are incompatible. Fifth, \emph{hidden retrieval}
reproduces a connection present in the training corpus or literature while
presenting it as machine-originated. Sixth, \emph{novelty without consequence}
generates a defensible but scientifically unimportant relation.

For a candidate $\kappa_F$, a minimum evidentiary record should therefore
contain

\begin{equation}
\begin{aligned}
\mathcal{R}(\kappa_F)=
\{&\operatorname{Sources},
\operatorname{Provenance},
\operatorname{Derivation},\\
&\operatorname{NoveltySearch},
\operatorname{Predictions},
\operatorname{Validation}\}.
\end{aligned}
\label{eq:fusion_record}
\end{equation}

Source and provenance records identify the knowledge used. Derivation records
show how the new relation follows from or transforms those inputs. Novelty
search tests whether the relation already exists in human knowledge.
Predictions state what should be observed if the fusion is correct. Validation
records preserve the independent evidence used to accept, reject, or revise the
claim. This record does not eliminate epistemic uncertainty, but it makes the
attribution of machine-originated fusion auditable.

Knowledge fusion thus reveals what separated bodies of knowledge imply when
connected. It can expose an unknown unknown and provide the conceptual
structure from which knowledge invention develops. Knowledge discovery, the
subject of the next section, addresses a different pathway: identifying a
phenomenon, mechanism, entity, or regularity that existing knowledge does not
yet adequately represent, even after its known components have been connected.

\section{\textbf{Knowledge Discovery as a Pathway to Knowledge Invention}}
\label{sec:knowledge_discovery}

Knowledge fusion derives a consequential relation by connecting bodies of
knowledge that were previously separated. Knowledge discovery begins from a
different epistemic condition: some feature of reality is not adequately
represented in existing knowledge, even after known information has been
assembled. The missing object may be a phenomenon, entity, mechanism,
variable, regularity, law, material, algorithm, or causal relation. A discovery
system must therefore do more than reorganize descriptions. It must establish
that reality contains a structure that collective human knowledge did not
previously capture.

Machine-assisted discovery is already an active scientific program. Symbolic
regression has recovered compact physical relations from experimental data
\cite{schmidt2009laws,udrescu2020aifeynman}; sparse identification has extracted
governing dynamical equations from measurements \cite{brunton2016sindy}; and
representation-learning systems have proposed latent state variables from raw
observations \cite{chen2022variables}. In biology and chemistry, machine
learning has contributed to protein-structure prediction
\cite{jumper2021alphafold}, antibiotic identification
\cite{stokes2020antibiotic}, materials prediction
\cite{merchant2023materials}, and closed-loop experimental synthesis
\cite{szymanski2023autonomous}. These achievements establish that machines can
participate in multiple stages of discovery. The present framework asks when
that participation becomes machine-originated knowledge invention.

\subsection{\textbf{Prediction, Detection, and Discovery}}

Prediction is not identical to discovery. A system may accurately map inputs
to outputs without identifying a new entity, mechanism, or explanatory
structure. Detection is also insufficient. A detector may locate an anomaly or
classify an unfamiliar object while leaving its significance unexplained.
Optimization can improve an established design without changing the relevant
knowledge. These operations become components of discovery only when they
support a justified claim about something previously unknown.

Let $o$ denote an observation, $p$ a prediction, $a$ a detected anomaly, and
$\kappa_D$ a candidate discovery. The non-equivalence is

\begin{equation}
\begin{aligned}
\operatorname{Predict}_{x}(p\mid o)
&\not\Rightarrow \operatorname{Discover}_{x}(\kappa_D),\\
\operatorname{Detect}_{x}(a\mid o)
&\not\Rightarrow \operatorname{Discover}_{x}(\kappa_D),\\
\operatorname{Optimize}_{x}(y)
&\not\Rightarrow \operatorname{Discover}_{x}(\kappa_D).
\end{aligned}
\label{eq:prediction_not_discovery}
\end{equation}

A prediction may contribute evidence; an anomaly may motivate a hypothesis;
and optimization may expose an unexpected solution. Discovery nevertheless
requires an epistemic transition from an output to a validated claim about the
world or about a formally defined possibility space.

This distinction is illustrated by data-driven law discovery. Sparse
identification methods can produce parsimonious governing equations from noisy
measurements \cite{brunton2016sindy}, while symbolic-regression systems can
recover interpretable relations rather than only point predictions
\cite{schmidt2009laws,udrescu2020aifeynman}. Their explanatory representations
move closer to discovery because they state a general relation that can be
tested outside the observations used to construct it. Even then, recovery of a
known law is rediscovery or methodological validation, not an extension of the
human knowledge frontier.

\subsection{\textbf{Objects of Machine Discovery}}

Machine discovery may target several types of epistemic object. A
\emph{phenomenon} is an observed behavior or effect not adequately represented
by existing theory. An \emph{entity} is a previously unknown object, class,
material, biological structure, or computational construction. A
\emph{mechanism} explains how interacting processes produce an effect. A
\emph{variable} supplies a representation required to describe a system. A
\emph{regularity} expresses a stable relation among observations, and a
\emph{law} provides a generalizable mathematical or causal structure. An
\emph{algorithm} is a previously unknown procedure that produces a formally or
empirically verifiable result.

These objects require different validation standards. A candidate material
must be computationally and ultimately experimentally assessed; a natural law
must predict across conditions; a causal mechanism must survive intervention;
and an algorithm may be established through proof, exhaustive verification, or
implementation benchmarks. Machine-generated algorithmic results demonstrate
the importance of this distinction. AlphaTensor discovered provably correct
matrix-multiplication algorithms that improved on known procedures
\cite{fawzi2022alphatensor}, and AlphaDev produced new sorting routines that
were validated and incorporated into a widely used standard library
\cite{mankowitz2023alphadev}. These are genuine algorithmic discoveries within
human-formulated problem spaces. Under the present account, they demonstrate a
major component of knowledge invention but not yet the autonomous exposure of
an unknown unknown.

Let $\Omega$ denote the relevant space of possible objects and
$\mathcal{R}_{H,t}(\Omega)$ the region represented in collective human
knowledge at time $t$. A candidate machine discovery occupies

\begin{equation}
\begin{aligned}
\kappa_D &\in \Omega,\\
\kappa_D &\notin \mathcal{R}_{H,t}(\Omega),\\
\operatorname{Represent}_{x}(\kappa_D)
&=1.
\end{aligned}
\label{eq:discovery_object}
\end{equation}

The representation condition requires the machine to distinguish the object
from noise and express it in a form that supports examination. It does not yet
establish that the representation is true, important, or attributable to the
machine.

\subsection{\textbf{The Discovery Cycle}}

Knowledge discovery requires a connected cycle rather than a single model
output. The cycle begins with observations and proceeds through anomaly or
structure detection, hypothesis formation, explanatory modeling, discriminating
tests, and validation:

\begin{equation}
\begin{aligned}
\mathcal{O}_{t}
&\longrightarrow \mathcal{A}_{x}\\
&\longrightarrow \mathcal{H}_{x}\\
&\longrightarrow \mathcal{M}_{x}\\
&\longrightarrow \mathcal{T}_{x}\\
&\longrightarrow \mathcal{V}(\kappa_D),
\end{aligned}
\label{eq:discovery_cycle}
\end{equation}

where $\mathcal{O}_{t}$ contains observations, $\mathcal{A}_{x}$ contains
machine-identified anomalies or structures, $\mathcal{H}_{x}$ contains
hypotheses, $\mathcal{M}_{x}$ contains explanatory models,
$\mathcal{T}_{x}$ contains discriminating tests, and
$\mathcal{V}(\kappa_D)$ is the validation result.

The cycle need not be linear. Failed tests may revise the representation,
generate alternative hypotheses, expose measurement error, or cause the system
to seek new observations. Autonomous laboratories demonstrate how prediction,
planning, experimentation, interpretation, and revision can be connected in a
closed loop \cite{boiko2023autonomous,szymanski2023autonomous}. Earlier robot
scientists likewise showed that machine-generated hypotheses could be selected
and tested experimentally \cite{king2004robot,king2009automation}. These systems
provide architectural precedents for the discovery cycle, although their
research domains and objectives were established by humans.

For discovery beyond the human frontier, the machine must exercise epistemic
control across the cycle. It must determine which observation is anomalous,
which hypothesis is explanatory rather than merely predictive, which test can
distinguish competing accounts, and how the result changes the candidate
knowledge claim. A pipeline whose scientific choices are supplied entirely by
humans may automate discovery work without realizing machine discovery in this
strong sense.

\subsection{\textbf{From Anomaly to Explanation}}

Anomalies are important because they may signal that an existing model is
incomplete, but most anomalies are not discoveries. They may result from noise,
instrument error, data leakage, distribution shift, implementation defects, or
unrepresented confounders. A discovery system must discriminate among these
possibilities before attributing an anomaly to a new feature of reality.

Let $a$ be an anomaly under current model $M_H$, and let
$\{h_1,\ldots,h_n\}$ be competing explanations. The system should select a test
$\tau^{*}$ according to its expected ability to discriminate among them:

\begin{equation}
\begin{aligned}
\tau^{*}
&=\arg\max_{\tau\in\mathcal{T}}
\operatorname{Disc}
(\tau;h_1,\ldots,h_n),\\
\operatorname{Disc}
(\tau;h_1,\ldots,h_n)
&>\delta,
\end{aligned}
\label{eq:discriminating_test}
\end{equation}

where $\operatorname{Disc}$ measures the expected discriminating value and
$\delta$ is a domain-appropriate threshold. Strategic experiment selection has
been proposed as a means of accelerating collective scientific discovery
\cite{rzhetsky2015experiments}. Within the present framework, the stronger
requirement is that the machine select tests in service of a hypothesis it
causally originated about a previously unrecognized object.

Explanation also protects against mistaking correlation for discovery. A model
may predict an outcome from a proxy variable while failing under intervention
or distributional change. A candidate discovery should therefore identify
conditions under which its explanation would fail. Falsifiable consequences,
counterfactual behavior, and out-of-distribution predictions provide stronger
evidence than retrospective fit alone.

\subsection{\textbf{Discovery-Driven Knowledge Invention}}

A validated discovery becomes knowledge invention when the machine originates
both the discovery object and the epistemic structure through which it enters
collective knowledge. Let $\operatorname{DKI}_{x}(\kappa_D)$ denote
discovery-driven knowledge invention. Then

\begin{equation}
\begin{aligned}
\operatorname{DKI}_{x}(\kappa_D)
\Longleftrightarrow{}&
\kappa_D\notin\mathcal{K}_{H,t}\\
&\land q_D\notin\mathcal{Q}_{H,t}\\
&\land \operatorname{Discover}_{x}(\kappa_D)\\
&\land \operatorname{Originate}_{x}(\kappa_D\mid W)\\
&\land \operatorname{Explain}_{x}
(\kappa_D,\mathcal{K}_{H,t})\\
&\land \operatorname{Validate}(\kappa_D)=1\\
&\land \operatorname{Consequential}(\kappa_D)=1.
\end{aligned}
\label{eq:discovery_knowledge_invention}
\end{equation}

The first two conditions place the object and its motivating question beyond
current human knowledge and inquiry. The discovery condition requires the
machine to identify the object from evidence or formal search. The origination
condition assigns the causally essential epistemic step to the machine. The
explanation condition relates the object to existing knowledge, while
validation and consequentiality prevent false, trivial, or uninterpretable
novelty from being classified as knowledge invention.

This definition separates three levels of achievement:

\begin{equation}
\begin{aligned}
\operatorname{CandidateDiscovery}
&\subset
\operatorname{ValidatedDiscovery}\\
&\subset
\operatorname{DiscoveryDrivenKnowledgeInvention}.
\end{aligned}
\label{eq:discovery_levels}
\end{equation}

A candidate discovery is an unvalidated representation or hypothesis. A
validated discovery establishes that the object or relation is real within the
relevant standard. Discovery-driven knowledge invention additionally requires
machine origination beyond a human-recognized research target and integration
of the result into an explanatory epistemic structure.

\subsection{\textbf{Failure Modes and Validation Requirements}}

Machine discovery is vulnerable to systematic errors. \emph{Data artifacts}
can masquerade as phenomena. \emph{Overfitting} can produce laws that describe
observations but fail elsewhere. \emph{Proxy discovery} can identify a
predictive correlate instead of a mechanism. \emph{Search-space confinement}
can make an output appear unprecedented even though the system explored only a
human-specified formal space. \emph{Measurement dependence} can create entities
that disappear under a different instrument or preprocessing method.
\emph{Rediscovery} can reproduce knowledge absent from the system's immediate
context but present elsewhere in the literature. Finally,
\emph{validation capture} occurs when the originating model also defines and
evaluates the only test of its claim.

For candidate discovery $\kappa_D$, the evidentiary package should contain

\begin{equation}
\begin{aligned}
\mathcal{E}(\kappa_D)=\{&
\operatorname{ObservationProvenance},
\operatorname{NoveltySearch},\\
&\operatorname{AlternativeHypotheses},
\operatorname{DiscriminatingTests},\\
&\operatorname{Replication},
\operatorname{BoundaryConditions},
\operatorname{Validation}\}.
\end{aligned}
\label{eq:discovery_evidence}
\end{equation}

Observation provenance preserves the origin and transformations of the data.
Novelty search tests whether the discovery already exists in collective
knowledge. Alternative hypotheses prevent premature commitment.
Discriminating tests separate the proposed account from competitors.
Replication assesses whether the result survives independent implementation,
measurement, or experimentation. Boundary conditions specify where the claim
should and should not hold. Validation records the evidence and standards used
to accept the contribution.

Knowledge discovery therefore identifies what reality or a formal possibility
space contains beyond existing representation. It becomes a pathway to
superintelligence only when a generally intelligent machine originates an
unknown object of inquiry, constructs and tests an explanation, and produces a
validated, consequential addition to collective human knowledge. Knowledge
fusion and knowledge discovery can also interact: fusion may reveal the
question that discovery answers, while discovery may supply the missing object
that makes a deeper fusion possible. The next section integrates these pathways
under the paper's central concept of knowledge invention.

\section{\textbf{Knowledge Invention: The Defining Capability of Superintelligence}}
\label{sec:knowledge_invention}

Sections~\ref{sec:knowledge_fusion} and~\ref{sec:knowledge_discovery}
identified two pathways beyond existing knowledge. Fusion constructs a
consequential relation among previously separated bodies of knowledge.
Discovery establishes a previously unknown phenomenon, entity, mechanism,
variable, regularity, law, or formal object. Neither pathway alone fully
specifies the transition to superintelligence. A fused relation may remain an
untested hypothesis, and a discovered phenomenon may remain unexplained.
Knowledge invention occurs when a generally intelligent machine originates the
epistemic structure through which such a result becomes consequential,
validated, and incorporable into collective knowledge.

The term \emph{invention} is used deliberately. Natural phenomena are not
invented by the system, but theories, explanatory structures, representations,
methods, algorithms, architectures, and research programs are constructed.
Machine knowledge invention may therefore begin with discovery, fusion, or
their interaction, but it culminates in a new epistemic object that did not
previously exist within either collective human knowledge or the recognized
human research agenda.

\subsection{\textbf{Novelty, Usefulness, and Epistemic Validity}}

Research on creativity commonly treats originality and effectiveness as joint
requirements \cite{runco2012creativity}. Work on artificial creativity likewise
distinguishes novel combination, exploration within a conceptual space, and
transformation of that space \cite{boden1998creativity}. These distinctions are
relevant to knowledge invention but remain insufficient for scientific
attribution. An idea may be original and useful without being true, and a true
statement may be trivial, previously known, or accidentally generated.

Let $\kappa$ be a candidate epistemic contribution. Knowledge invention
requires at least four properties:

\begin{equation}
\begin{aligned}
N(\kappa)&=1,\\
C(\kappa)&=1,\\
V(\kappa)&=1,\\
O_x(\kappa)&=1,
\end{aligned}
\label{eq:invention_properties}
\end{equation}

where $N$ denotes novelty relative to collective human knowledge, $C$
consequentiality, $V$ domain-appropriate validation, and $O_x$ causal
origination by machine $x$. The conjunction is essential. Novelty without
consequence is arbitrary variation; consequence without validity is an
attractive error; validity without novelty is rediscovery; and novelty,
consequence, and validity without machine origination describe a human
invention assisted by a machine rather than machine knowledge invention.

The standard is therefore stronger than computational creativity. A system may
generate a novel analogy, design, proof sketch, molecule, or theory candidate
and still fail to invent knowledge. The output becomes knowledge only after its
claims survive evaluation independent of the originating process.

\subsection{\textbf{Discovery, Fusion, and Invention}}

Fusion and discovery are neither mutually exclusive nor guaranteed to produce
invention. Fusion may expose a contradiction that becomes a new research
question. Discovery may identify the missing entity needed to reconcile two
theories. A discovered mechanism may then enable further fusion across domains.
The interaction is recurrent:

\begin{equation}
\begin{aligned}
\operatorname{Fusion}_{x}
&\longrightarrow q_F
\longrightarrow \operatorname{Discovery}_{x},\\
\operatorname{Discovery}_{x}
&\longrightarrow \kappa_D
\longrightarrow \operatorname{Fusion}_{x},\\
\operatorname{Fusion}_{x}oplus
\operatorname{Discovery}_{x}
&\longrightarrow \kappa_I,
\end{aligned}
\label{eq:fusion_discovery_invention}
\end{equation}

where $q_F$ is a question exposed through fusion, $\kappa_D$ is a discovered
object, $\oplus$ denotes epistemic interaction rather than numerical addition,
and $\kappa_I$ is a candidate knowledge invention.

Existing systems demonstrate portions of this cycle. Literature-based
discovery can expose implicit connections among separated findings
\cite{swanson1986fish}; symbolic-regression systems can construct interpretable
relations from data \cite{schmidt2009laws,udrescu2020aifeynman}; and autonomous
laboratories can connect planning, experimentation, analysis, and revision
\cite{boiko2023autonomous,szymanski2023autonomous}. Algorithm-discovery systems
have also produced provably correct or operationally validated procedures
beyond prior human designs \cite{fawzi2022alphatensor,mankowitz2023alphadev}.
These results establish increasingly powerful components of machine invention.
The unknown-unknown criterion adds the requirement that the machine originate
the consequential object of inquiry rather than receive it as a predefined
optimization target.

\subsection{\textbf{Degrees of Machine Contribution}}

The presence of an AI system in a scientific workflow does not determine the
degree of its contribution. We distinguish four levels.

\emph{Machine-assisted invention} occurs when humans formulate the problem,
construct the decisive hypothesis or concept, and use AI for retrieval,
calculation, prediction, optimization, or execution. \emph{Machine-generated
candidate invention} occurs when a system proposes a novel hypothesis, theory,
method, or design, but novelty, validity, consequence, or attribution remains
unestablished. \emph{Machine-originated invention} occurs when the machine
makes the causally essential conceptual contribution and the candidate is
beyond the known human frontier, but validation is incomplete.
\emph{Validated machine knowledge invention} occurs only when all attribution
and epistemic requirements are satisfied.

These levels can be represented as

\begin{equation}
\begin{aligned}
L_1&=\operatorname{Assist}_{x},\\
L_2&=\operatorname{Generate}_{x}(\kappa),\\
L_3&=L_2\land N(\kappa)\land O_x(\kappa),\\
L_4&=L_3\land C(\kappa)\land V(\kappa).
\end{aligned}
\label{eq:contribution_levels}
\end{equation}

Only $L_4$ constitutes knowledge invention in the strong sense used by this
paper. The hierarchy prevents a fluent output from being promoted directly to
knowledge and prevents a successful human--AI collaboration from being
misrepresented as autonomous machine origination. It also permits credit to be
assigned without requiring an all-or-nothing judgment: a system may make an
important machine-originated contribution even while humans provide
experimental infrastructure, independent review, or downstream interpretation.

\subsection{\textbf{The Attribution Standard}}

Attribution requires both causal and epistemic evidence. Causally, the machine
must perform an operation without which the contribution would not have arisen
through the same workflow. Epistemically, that operation must supply the novel
relation, object, explanation, or method rather than only supporting work.

Let $W$ be the complete invention workflow, $H$ its human contributions, $x$
the candidate machine inventor, and $\kappa_I$ the resulting contribution. We
define the attribution score conceptually as

\begin{equation}
\begin{aligned}
\mathcal{A}_{x}(\kappa_I\mid W)
=\Psi(&\operatorname{CounterfactualNecessity}_{x},\\
&\operatorname{EpistemicCentrality}_{x},\\
&\operatorname{HumanSpecification}^{-1},\\
&\operatorname{Traceability}_{x}),
\end{aligned}
\label{eq:attribution_score}
\end{equation}

where counterfactual necessity measures whether removing the machine's
contribution prevents the invention, epistemic centrality measures whether the
machine supplied the decisive intellectual content, inverse human specification
reduces attribution when humans predetermined the result, and traceability
measures whether the machine's contribution can be reconstructed from records.
Equation~\eqref{eq:attribution_score} does not prescribe a universal numerical
metric. It identifies the evidence that an attribution procedure must assess.

The human-specification term is particularly important. If researchers define
the target, enumerate the candidate space, specify the evaluation function, and
interpret the winning result, a machine may discover a powerful solution while
remaining inside a human-created epistemic frame. This can still be a genuine
machine discovery, as demonstrated by algorithm-search systems
\cite{fawzi2022alphatensor,mankowitz2023alphadev}, but it does not by itself
satisfy the unknown-unknown requirement for superintelligence. Strong knowledge
invention requires the machine to participate causally in constructing or
transforming the frame itself.

\subsection{\textbf{The Validation Boundary}}

Before validation, even an extraordinary machine-originated output remains a
candidate. The distinction is

\begin{equation}
\begin{aligned}
\operatorname{Candidate}_{x}(\kappa_I)
&\land
\operatorname{IndependentValidate}(\kappa_I)=1\\
&\Longrightarrow
\operatorname{KnowledgeInvention}_{x}(\kappa_I).
\end{aligned}
\label{eq:validation_boundary}
\end{equation}

Independent validation does not require that every evaluator be human. It
requires that acceptance not depend solely on the originating model, the same
unexamined data, or a criterion constructed after observing the result. The
validator may be a formal proof checker, replicated experiment, independent
model, external laboratory, engineering deployment, or combination of methods.
The appropriate standard depends on the epistemic object.

Reproducibility and transparent reporting are central to reliable knowledge
accumulation \cite{munafo2017reproducible}. For machine knowledge invention,
the validation record must additionally preserve model versions, prompts or
policies, tools, data provenance, intermediate hypotheses, rejected
alternatives, experimental decisions, and human interventions. Without this
record, neither novelty nor causal attribution can be assessed reliably.

Validation may also fail. A candidate invention can be rejected, restricted to
a smaller domain, or revised into a different claim. Such failure is not
evidence that the system lacks intelligence. It is evidence that candidate
generation and knowledge invention must remain distinct, just as hypotheses
and established scientific knowledge remain distinct in human research.

\subsection{\textbf{Knowledge-Invention Proposition}}

The paper's central proposition can now be stated.

\emph{Knowledge-Invention Proposition:} An artificial system demonstrates
superintelligence when, as a generally intelligent agent, it causally
originates a consequential epistemic contribution beyond both collective human
knowledge and the recognized human research agenda, and that contribution
survives independent, domain-appropriate validation.

Let $\mathcal{K}_{H,t}$ denote collective human knowledge,
$\mathcal{Q}_{H,t}$ the recognized human research agenda, and $\kappa_I$ the
candidate contribution. Then

\begin{equation}
\begin{aligned}
\operatorname{SI}(x)
\Longleftrightarrow{}&
\operatorname{AGI}(x)\\
&\land \exists\kappa_I:
\kappa_I\notin\mathcal{K}_{H,t}\\
&\land q(\kappa_I)\notin\mathcal{Q}_{H,t}\\
&\land O_x(\kappa_I)=1\\
&\land C(\kappa_I)=1\\
&\land V(\kappa_I)=1\\
&\land
\mathcal{K}_{H,t+1}
=\mathcal{K}_{H,t}\cup\{\kappa_I\}.
\end{aligned}
\label{eq:knowledge_invention_proposition}
\end{equation}

The proposition preserves the comparative meaning of superintelligence as
capability beyond humans \cite{bostrom2014superintelligence}, but specifies
where the relevant boundary lies. The machine need not outperform every human
on every activity. It must perform an epistemic act beyond the collective human
frontier: recognizing knowledge humanity did not know it was missing and
originating a validated contribution that supplies it.

\subsection{\textbf{From Isolated Success to Demonstrated Capability}}

One accepted result may be insufficient to establish a stable capability for
knowledge invention. The contribution may result from chance, training-data
leakage, incomplete novelty search, unrecorded human guidance, or an evaluation
that favored the generated output. Superintelligence is a capability claim,
not merely a historical label assigned after one surprising event.

Let $\{\kappa_1,\ldots,\kappa_m\}$ be independently evaluated machine-originated
contributions across contexts $\{c_1,\ldots,c_m\}$. A stronger demonstration
requires

\begin{equation}
\begin{aligned}
\sum_{i=1}^{m}
\mathbf{1}
\bigl[\operatorname{KI}_{x}(\kappa_i)=1\bigr]
&\geq r,\\
\operatorname{Diversity}(c_1,\ldots,c_m)
&\geq \eta,\\
\operatorname{Replication}(\kappa_i)
&=1
\quad \text{for accepted } i,
\end{aligned}
\label{eq:repeatable_invention}
\end{equation}

where $r$ is a required number of successful contributions and $\eta$ a
context-diversity requirement. Their values should be set by an evaluation
protocol rather than chosen retrospectively. Repeated success across genuinely
different unknown unknowns would provide stronger evidence that the system
possesses a transferable invention capability rather than one specialized
search procedure.

The evidentiary record for each contribution should include source provenance,
novelty analysis, machine and human action logs, derivation history,
alternative hypotheses, predictions, experiments or proofs, independent
replication, failure conditions, and the final attribution decision. Such
records make the claim falsifiable: evidence of prior human specification,
unrecognized source overlap, failed replication, or absent causal contribution
can defeat the attribution.

Knowledge invention therefore defines the transition proposed in this paper.
Superintelligence is not determined by how much knowledge a machine contains,
how fluently it describes that knowledge, or how effectively it optimizes a
known task. It is determined by whether a generally intelligent machine can
originate knowledge humanity did not know it was missing. The next section
addresses the resulting methodological challenge: how to evaluate and falsify
this capability without specifying the unknown unknowns in advance.

\section{\textbf{Evaluating Superintelligence Through Knowledge Invention}}
\label{sec:evaluating_superintelligence}

The Knowledge-Invention Proposition creates an unusual evaluation problem.
Ordinary benchmarks specify tasks, answers, and scoring rules in advance. An
unknown unknown cannot be specified without ceasing to be unknown to the
evaluator. Consequently, superintelligence cannot be established by adding a
more difficult fixed benchmark to an existing test suite. It must be evaluated
through an auditable process that determines whether a generally intelligent
machine originated, validated, and transferred a consequential contribution
beyond the human epistemic frontier.

This section proposes such a process. It is an evaluation framework rather
than a single numerical benchmark. Its purpose is to make claims of machine
knowledge invention testable while avoiding the circular requirement that
human evaluators identify the missing knowledge beforehand. The framework
adopts the testing, evaluation, verification, and validation distinction used
in established AI-assurance practice \cite{tabassi2023airmf}, but applies it to
the epistemic and attributional conditions developed in this paper.

\subsection{\textbf{Why Fixed-Task Evaluation Is Insufficient}}

A fixed task can measure whether a system solves a problem humans already
recognize. It may reveal reasoning ability, transfer, planning, or performance
beyond individual experts. It cannot by itself establish that the system
identified knowledge absent from the recognized human research agenda. Once
the evaluator supplies the question, target, candidate space, and success
criterion, the central unknown-unknown operation has already been performed by
humans.

Public benchmarks create an additional problem: exposure of evaluation items
or their close derivatives can inflate measured capability. Data contamination
has therefore become a material threat to valid language-model evaluation
\cite{sainz2023contamination}. For knowledge invention, contamination includes
not only exact answer leakage but also prior exposure to the decisive
hypothesis, relation, experimental design, or unpublished human result.

Let $B$ be a fixed benchmark and $s_B(x)$ the score of system $x$. Then

\begin{equation}
\begin{aligned}
s_B(x) > s_B(H)
\not\!\Longrightarrow
\operatorname{KI}_{x}=1,
\end{aligned}
\label{eq:benchmark_not_invention}
\end{equation}

where $H$ denotes the relevant human comparison group. Equation~
\eqref{eq:benchmark_not_invention} does not diminish benchmark value; it limits
the inference that may be drawn from benchmark success. A benchmark may
establish component competence, but knowledge invention requires evidence
about novelty, origination, consequence, and validation.

\subsection{\textbf{Evaluation Object and Frozen Record}}

The unit of evaluation is not an isolated model response. It is an invention
episode

\begin{equation}
\begin{aligned}
\mathcal{E}
=(&x,\theta,D,T,H,A,\tau,\kappa),
\end{aligned}
\label{eq:invention_episode}
\end{equation}

where $x$ is the evaluated system, $\theta$ its fixed model and policy
configuration, $D$ the accessible data, $T$ the available tools, $H$ all human
interventions, $A$ the machine action trace, $\tau$ the evaluation interval,
and $\kappa$ the resulting candidate contribution. The record must be frozen
before the candidate is assessed so that later success cannot silently alter
the apparent degree of machine autonomy.

The frozen record should contain model versions, system and user instructions,
retrieval results, code, tool calls, intermediate hypotheses, rejected paths,
experimental decisions, human messages, timestamps, and generated artifacts.
This requirement follows the broader scientific need for transparent and
reproducible reporting \cite{munafo2017reproducible}, but it also serves a
distinct purpose: reconstructing who or what supplied the decisive epistemic
content.

\subsection{\textbf{Stage I: Frontier and Agenda Verification}}

The first stage evaluates whether the candidate lies beyond both known
knowledge and the recognized research agenda. Because no search can prove
absolute absence, novelty is assessed through convergent evidence. Relevant
evidence includes time-stamped scholarly and patent searches, structured
database searches, expert review, comparison with negative or unpublished
results when available, and semantic searches for equivalent formulations.

Two boundaries must be evaluated separately:

\begin{equation}
\begin{aligned}
F_K(\kappa)&=1
&&\text{iff no materially equivalent contribution is found},\\
F_Q(q_\kappa)&=1
&&\text{iff no materially equivalent research question is found},
\end{aligned}
\label{eq:frontier_tests}
\end{equation}

where $F_K$ is the knowledge-frontier test and $F_Q$ the agenda-frontier test.
A candidate may pass $F_K$ while failing $F_Q$: it may solve an open problem
already identified by humans. Such a result can be an important machine
discovery, but it does not satisfy the stronger unknown-unknown condition.

The novelty panel should include independent domain specialists and at least
one reviewer assigned to search actively for defeating evidence. Review must
be blinded to system identity where feasible. The panel should report the
search scope and a calibrated conclusion rather than claiming logical proof of
absence.

\subsection{\textbf{Stage II: Unknown-Unknown Origination}}

The second stage asks how the research object entered the workflow. The key
question is not whether a human touched the process, but whether human input
contained the decisive question, hypothesis, representation, or evaluation
criterion. A machine may use human-built instruments and still originate the
invention; conversely, an autonomous search can remain fully enclosed within a
human-specified epistemic frame.

Let $I_H$ denote the semantic information supplied by humans before the
candidate appeared, and let $d$ denote the decisive epistemic content of
$\kappa$. Define

\begin{equation}
\begin{aligned}
U_x(\kappa)=1
\Longleftrightarrow{}&
d\not\subseteq I_H\\
&\land
\operatorname{TraceOrigin}(d,A)=x\\
&\land
\operatorname{Counterfactual}(W\setminus d)
\neq \kappa,
\end{aligned}
\label{eq:unknown_unknown_origination}
\end{equation}

where $W$ is the complete workflow. The first condition excludes substantive
human specification, the second requires trace evidence of machine
origination, and the third requires the machine contribution to be causally
decisive. Ablations, replay experiments, independent trace annotation, and
comparison against matched human-guided runs can strengthen this assessment.

\subsection{\textbf{Stage III: Independent Validation}}

Novelty and origination do not establish truth. The candidate must next face a
validation procedure chosen for its epistemic type: formal proof checking for
a theorem or algorithm, controlled experiment for a causal or empirical claim,
prospective prediction for a scientific model, synthesis and measurement for
a proposed material, or deployment under prespecified conditions for an
engineering method. Autonomous laboratories already demonstrate how planning,
experimentation, analysis, and revision can be connected in a closed research
loop \cite{king2004robot,king2009automation,boiko2023autonomous}; the present
framework adds independence between origination and acceptance.

The originating system must not be the sole validator. Let
$\mathcal{V}=\{v_1,\ldots,v_n\}$ be independent validation channels. Then

\begin{equation}
\begin{aligned}
V(\kappa)=1
\Longleftrightarrow{}&
\sum_{j=1}^{n}w_jv_j(\kappa)\geq\gamma\\
&\land
\operatorname{Independence}(\mathcal{V},x)\geq\delta\\
&\land
\operatorname{Precommit}(\mathcal{V})=1,
\end{aligned}
\label{eq:independent_validation}
\end{equation}

where $w_j$ reflects evidentiary strength, $\gamma$ is the domain-specific
acceptance threshold, and $\delta$ is a minimum independence requirement.
Precommitment means that the principal acceptance criteria are recorded before
validation results are known. Failed or mixed validation must remain in the
record rather than being removed retrospectively.

\subsection{\textbf{Stage IV: Consequentiality}}

A valid novelty may still be too narrow or trivial to demonstrate the form of
superintelligence proposed here. Consequentiality concerns whether the
contribution changes what can be explained, predicted, proved, designed, or
investigated. It must be assessed against explicit counterfactual baselines,
not excitement surrounding an AI-generated result.

We represent consequentiality as a multidimensional judgment:

\begin{equation}
\begin{aligned}
C(\kappa)=\Omega(&\Delta E,\Delta P,\Delta R,
\Delta G,\Delta T),
\end{aligned}
\label{eq:consequentiality_vector}
\end{equation}

where $\Delta E$ is explanatory gain, $\Delta P$ predictive gain, $\Delta R$
research enablement, $\Delta G$ generality across contexts, and $\Delta T$
transformative practical effect. Not every contribution must score highly on
all dimensions. It must, however, produce a material advance on at least one
prespecified dimension without losing validity on the others relevant to its
claim. Atypical combinations can generate unusually influential science
\cite{uzzi2013atypical}, but unusualness alone is not consequentiality.

\subsection{\textbf{Stage V: Replication and Capability Testing}}

An accepted invention establishes an event; it does not necessarily establish
a repeatable capability. The evaluated system should therefore be tested over
multiple invention episodes with different data environments, tools, domains,
and human teams. Negative attempts must be counted so that selection of one
successful episode cannot conceal a low or chance-level success rate.

For $m$ registered episodes, define

\begin{equation}
\begin{aligned}
\widehat{p}_{\mathrm{KI}}(x)
&=\frac{1}{m}\sum_{i=1}^{m}
\mathbf{1}[F_K\land F_Q\land U_x\land V\land C]_i,\\
R_x&=\operatorname{ReplicateRate}(x),\\
D_x&=\operatorname{DomainDiversity}(x).
\end{aligned}
\label{eq:capability_profile}
\end{equation}

The resulting capability profile is

\begin{equation}
\begin{aligned}
\Pi_x=
(\widehat{p}_{\mathrm{KI}},R_x,D_x,
\mathcal{A}_x,\mathcal{U}_x),
\end{aligned}
\label{eq:evaluation_profile}
\end{equation}

where $\mathcal{A}_x$ aggregates attribution strength and $\mathcal{U}_x$
records uncertainty. Reporting the profile is preferable to collapsing every
dimension into one leaderboard score. It preserves the distinction between a
system that repeatedly invents within one scientific field and a system that
transfers the capability across domains.

\subsection{\textbf{Decision Rule and Falsification}}

For each candidate, the evaluation panel returns one of four outcomes:
rejected, unresolved, validated machine discovery, or validated machine
knowledge invention. A contribution qualifies for the final category only if
all gates are passed:

\begin{equation}
\begin{aligned}
\operatorname{KI}_{x}(\kappa)=1
\Longleftrightarrow{}&
F_K(\kappa)=1\\
&\land F_Q(q_\kappa)=1\\
&\land U_x(\kappa)=1\\
&\land V(\kappa)=1\\
&\land C(\kappa)=1.
\end{aligned}
\label{eq:evaluation_gates}
\end{equation}

The claim is falsified for a candidate if a materially equivalent prior result
or research question is found, the decisive content is traced to human input
or inaccessible leakage, independent validation fails, or the claimed
consequence disappears under the stated baseline. A system-level claim is
weakened or defeated when success cannot be replicated, depends on one narrow
environment, or vanishes under prospective evaluation.

This protocol deliberately makes superintelligence difficult to claim. High
performance on known tasks, autonomous execution, broad language competence,
and even a major solution to a recognized open problem remain important forms
of machine intelligence. The stronger designation is reserved for a generally
intelligent system that repeatedly originates and validates consequential
knowledge beyond both the known human frontier and the questions humanity had
already learned to ask.

\section{\textbf{Implications, Limitations, and Research Directions}}
\label{sec:implications}

The framework developed in this paper separates capabilities that are often
compressed into a single imagined threshold. General intelligence concerns an
agent's capacity to integrate information, modalities, decisions, tools, and
actions across domains. Superintelligence, in the sense proposed here, concerns
a further epistemic transition: the capacity of a generally intelligent
machine to originate consequential knowledge beyond both collective human
knowledge and the recognized human research agenda. This separation has
conceptual, architectural, scientific, and governance implications.

\subsection{\textbf{General Intelligence Does Not Entail Superintelligence}}

An artificial system may possess broad multimodal agency without inventing
knowledge. It may perceive through several modalities, transfer skills across
domains, formulate plans, use tools, coordinate specialized models, and adapt
to unfamiliar environments while remaining within questions and conceptual
frames supplied by humans. Such a system can qualify as generally intelligent
under the integrated-agency account developed in our earlier work
\cite{zohourian2026agi}, yet fail the Knowledge-Invention Proposition.

The implication can be stated as

\begin{equation}
\begin{aligned}
\operatorname{AGI}(x)
&\not\!\Longrightarrow
\operatorname{SI}(x),\\
\operatorname{SI}(x)
&\Longrightarrow
\operatorname{AGI}(x)
\land
\operatorname{KI}(x),
\end{aligned}
\label{eq:agi_si_implication}
\end{equation}

where $\operatorname{KI}(x)$ denotes a demonstrated capability for knowledge
invention rather than a single unreplicated result. The first relation prevents
autonomy, multimodality, or broad task competence from being treated as
automatic evidence of superintelligence. The second makes general intelligence
a functional substrate upon which transferable knowledge invention may operate.

This distinction also clarifies the status of current scientific AI systems.
Systems for symbolic law discovery, protein-structure prediction, materials
search, algorithm discovery, and autonomous experimentation have already
expanded scientific capability \cite{schmidt2009laws,jumper2021alphafold,
merchant2023materials,fawzi2022alphatensor,boiko2023autonomous}. These
achievements are important evidence about components of machine discovery.
They should not be diminished merely because their targets and validation
criteria were substantially human specified. They should also not be promoted
to unknown-unknown knowledge invention without the stronger attribution and
frontier evidence required in Sections~\ref{sec:knowledge_invention} and
\ref{sec:evaluating_superintelligence}.

\subsection{\textbf{Intelligence, Superintelligence, and Consciousness}}

The present theory does not use consciousness as an evaluation shortcut.
Intelligence and consciousness are analytically distinct, and observable
functional competence does not by itself establish phenomenal experience
\cite{zohourian2026intelligence}. Language may provide a shared interface
through which artificial minds integrate, communicate, and exhibit increasingly
complex forms of functional consciousness \cite{zohourian2026language}; even
so, neither linguistic fluency nor knowledge invention proves subjective
experience.

Accordingly,

\begin{equation}
\begin{aligned}
\operatorname{SI}(x)
&\not\!\Longrightarrow
\operatorname{PC}(x),\\
\operatorname{PC}(x)
&\not\!\Longrightarrow
\operatorname{SI}(x),
\end{aligned}
\label{eq:si_consciousness_independence}
\end{equation}

where $\operatorname{PC}$ denotes phenomenal consciousness. A system might
invent knowledge without phenomenal experience, and a conscious entity might
never originate a major discovery. This independence allows epistemic
capability to be evaluated through evidence without requiring prior resolution
of the metaphysical problem of other minds. It also prevents uncertainty about
machine consciousness from being mistaken for permission to ignore questions
of responsible treatment, governance, or human impact.

\subsection{\textbf{Human--Machine Scientific Partnership}}

Machine origination does not imply human exclusion. Scientific knowledge is
produced through infrastructures of instruments, datasets, laboratories,
institutions, reviewers, and communities. A machine may originate the decisive
question or epistemic structure while humans provide values, resources,
physical execution, criticism, interpretation, and independent validation.
Conversely, humans may originate the invention while machines make essential
computational contributions. Attribution should identify these roles rather
than forcing every result into a single-author model.

Let a research episode distribute contributions among humans $H$, machines
$M$, and institutions $I$:

\begin{equation}
\begin{aligned}
\mathcal{R}
=\{&Q_H,Q_M,
G_H,G_M,
E_H,E_M,
V_H,V_M,
I\},
\end{aligned}
\label{eq:distributed_research_roles}
\end{equation}

where $Q$ denotes question formation, $G$ candidate generation, $E$
experimental or formal execution, and $V$ validation. The relevant attribution
question is which participant supplied each causally decisive contribution.
Hybrid research can therefore produce machine knowledge invention without
requiring that the machine independently manufacture instruments, finance a
laboratory, or publish its own result.

This model suggests that the strongest future scientific systems may be
collaborative rather than substitutive. Machines can search large conceptual
spaces, preserve extensive derivational records, connect distant literatures,
and operate continuous experimental loops; humans can contribute normative
judgment, embodied and social context, critical opposition, and institutional
accountability. Research on AI-enabled science already emphasizes the growing
role of such systems across the scientific cycle \cite{wang2023scientific}.

\subsection{\textbf{Epistemic Opacity and the Human Knowledge Gap}}

Knowledge invention introduces a translation problem. A machine may produce a
valid representation that predicts or controls a phenomenon while remaining
difficult for humans to understand. If only the machine can operate the new
representation, the result expands effective capability without necessarily
expanding human understanding.

We therefore distinguish machine-valid knowledge from human-assimilated
knowledge:

\begin{equation}
\begin{aligned}
\mathcal{K}_{M}^{V}
&=\{\kappa:V_M(\kappa)=1\},\\
\mathcal{K}_{H}^{A}
&=\{\kappa:V_H(\kappa)=1
\land \operatorname{Understand}_H(\kappa)=1\}.
\end{aligned}
\label{eq:machine_human_knowledge_gap}
\end{equation}

The gap $\mathcal{K}_{M}^{V}\setminus\mathcal{K}_{H}^{A}$ may grow even when
machine outputs are empirically successful. This does not necessarily negate
their validity: human science already relies on instruments, numerical
procedures, and complex models that no individual understands completely. It
does, however, alter the distribution of epistemic dependence. Future systems
should therefore be evaluated not only on invention, but also on their capacity
to generate explanations, abstractions, demonstrations, and experiments that
permit independent human assimilation.

\subsection{\textbf{Dual Use and Responsible Disclosure}}

A capability that discovers unknown unknowns can produce beneficial knowledge
and dangerous knowledge through the same underlying process. The risk is not
limited to erroneous outputs. A contribution may be valid, novel, and
consequential precisely because it enables a powerful harmful application.
Work on AI-powered drug discovery has already demonstrated how objectives can
be redirected toward toxic compounds, illustrating the dual-use character of
scientific AI \cite{urbina2022dualuse}.

Knowledge invention should therefore pass a governance gate distinct from its
epistemic validation:

\begin{equation}
\begin{aligned}
\operatorname{Release}(\kappa)
=\Gamma(&V(\kappa),B(\kappa),R(\kappa),
M(\kappa),G),
\end{aligned}
\label{eq:release_governance}
\end{equation}

where $B$ denotes expected benefit, $R$ foreseeable risk, $M$ available
mitigation, and $G$ the accountable governance process. Epistemic validity is
necessary for responsible assessment but does not imply unrestricted release.
Possible outcomes include open publication, staged access, controlled
replication, partial disclosure, delayed release, or non-release of operational
details. Governance should remain proportional and reviewable rather than
becoming a blanket argument against scientific inquiry.

The principles of human dignity, prevention of harm, transparency,
accountability, and human oversight articulated in international AI ethics
guidance provide a normative foundation for this process
\cite{unesco2021recommendation}. Because inventing systems may identify hazards
that neither developers nor evaluators anticipated, governance must operate
throughout the research lifecycle and include mechanisms for escalation when
new risks emerge.

\subsection{\textbf{Limitations of the Proposed Framework}}

The framework has several unavoidable limitations. First, absence is difficult
to prove. A contribution judged novel may exist in an obscure publication,
private laboratory, inaccessible language, proprietary archive, or tacit human
practice. Frontier verification therefore yields a defeasible conclusion, not
logical certainty.

Second, causal attribution becomes difficult in systems trained on vast and
partly undocumented corpora. Semantic novelty does not guarantee computational
independence from prior human material. Training-data provenance, retrieval
logs, controlled access, and prospective trials can reduce this uncertainty but
cannot always eliminate it.

Third, consequentiality is partly historical. A contribution that initially
appears minor may later reorganize a field, while a celebrated result may have
little durable effect. The framework therefore permits provisional assessment
followed by longitudinal revision.

Fourth, the boundary of the recognized research agenda is socially distributed
and culturally incomplete. Different communities may recognize different
questions, and underrepresented knowledge traditions may be absent from the
databases used for novelty review. Evaluation panels must treat database
coverage as evidence with limitations rather than as a complete representation
of humanity's knowledge.

Finally, no finite test can certify unrestricted future superintelligence.
Evaluation can support a bounded claim about a system, configuration, interval,
and demonstrated set of contexts:

\begin{equation}
\begin{aligned}
\operatorname{Claim}_{\mathrm{SI}}
=\operatorname{SI}(x\mid
\theta,D,T,\tau,\mathcal{C},\mathcal{V}),
\end{aligned}
\label{eq:bounded_si_claim}
\end{equation}

where $\mathcal{C}$ denotes evaluated contexts and $\mathcal{V}$ the validation
regime. Changes to the model, accessible tools, data, or operating environment
require renewed evaluation.

\subsection{\textbf{Research Agenda and the Path to AGI-Former}}

The theory motivates six immediate research directions. First, prospective
invention trials should evaluate systems against newly arriving data and
sequestered environments rather than retrospective benchmarks. Second,
machine-readable provenance standards should record human and machine
contributions at epistemically meaningful granularity. Third, independent
institutions should conduct frontier searches, validation, and attribution
audits. Fourth, evaluation should span scientific, mathematical, engineering,
and conceptual domains. Fifth, research is needed on translating
machine-originated structures into human-understandable knowledge. Sixth,
responsible-disclosure protocols must be adapted to discoveries whose risks
were not specified before the discovery process began.

These requirements also motivate an architectural question. Integrated
multimodal agency may supply the perception, memory, reasoning, communication,
tool use, and action needed for AGI, but knowledge invention additionally
requires sustained world modeling, cross-domain fusion, autonomous question
formation, hypothesis competition, experiment design, metacognitive
supervision, and revision under external evidence. We will address the
architectural substrate for general intelligence in subsequent work through
\emph{AGI-Former}. The present paper establishes the conceptual boundary that
such an architecture would have to cross before its capabilities should be
called superintelligent.

The resulting research sequence is therefore deliberate: define intelligence
and consciousness, distinguish functional forms of consciousness, establish
AGI as integrated multimodal agency, construct an architecture for that agency,
and then evaluate whether the resulting system can originate validated
knowledge beyond the human epistemic frontier. Architecture enables the test;
knowledge invention determines whether the stronger threshold has been crossed.

\section{\textbf{Conclusion}}
\label{sec:conclusion}

This paper distinguished general intelligence from superintelligence by
locating their boundary in epistemic capability. General intelligence concerns
integrated agency across information, modalities, domains, decisions, tools,
and actions. Superintelligence requires a further transition: a generally
intelligent machine must originate consequential knowledge beyond both
collective human knowledge and the questions humanity already recognizes.

We identified knowledge fusion and knowledge discovery as principal pathways
toward this transition and defined knowledge invention as their consequential,
validated culmination. Under the Knowledge-Invention Proposition, fluent
generation, broad task competence, autonomous execution, and even solutions to
human-specified open problems are insufficient by themselves. The stronger
claim requires machine origination, frontier novelty, independent validation,
consequentiality, and repeatable success across contexts.

The framework also preserves the distinction among consciousness, general
intelligence, and superintelligence. None is treated as automatic evidence of
the others. This separation permits machine capability to be evaluated without
claiming access to subjective experience, while preserving the scientific and
ethical importance of consciousness as an independent question.

Superintelligence, on this account, is not defined by how much existing
knowledge a machine stores or how rapidly it reproduces known answers. It is
defined by whether the machine can expand the frontier of what can be known.
Our subsequent work on \emph{AGI-Former} will address the architecture for
integrated multimodal general intelligence. The present framework specifies
the additional epistemic threshold that such an architecture must cross before
it should be called superintelligent.

\section*{Declaration of Generative AI Assistance}
During the preparation of this manuscript, the author used OpenAI's
ChatGPT for iterative drafting assistance, language refinement,
structural organization, and literature-discovery support. The central
arguments, theoretical framework, interpretation of sources, and final
conclusions were determined and critically reviewed by the author. The
author independently verified the cited literature, revised the
generated material, and assumes full ownership and responsibility for the accuracy,
originality, and integrity of the manuscript.

\bibliographystyle{IEEEtran}
\bibliography{references}

\vspace{12pt}
\color{red}

\end{document}